\section{PROJECT BACKGROUND}

\subsection{Background}

\begin{frame}[t]{Background: KFCProcedure Algorithm Workflow}
    \footnotesize
    \centering

    {\scriptsize
        \textbf{Clusterwise learning framework based on KFCProcedure}
        \cite{has2021kfc}
    }

    \vspace{0.18cm}

    \resizebox{0.93\linewidth}{!}{
        \begin{tikzpicture}[
                font=\scriptsize,
                box/.style={
                        draw=blue!60,
                        fill=lightblue,
                        rounded corners=4pt,
                        minimum width=1.35cm,
                        minimum height=0.65cm,
                        align=center
                    },
                widebox/.style={
                        draw=blue!60,
                        fill=lightblue,
                        rounded corners=4pt,
                        minimum width=2.45cm,
                        minimum height=0.65cm,
                        align=center
                    },
                aggbox/.style={
                        draw=blue!60,
                        fill=green!12,
                        rounded corners=4pt,
                        minimum width=3.9cm,
                        minimum height=0.8cm,
                        align=center
                    },
                note/.style={
                        anchor=west,
                        align=left,
                        text width=2.75cm,
                        font=\scriptsize
                    },
                arrow/.style={-Latex, thick},
                dashline/.style={dash pattern=on 5pt off 4pt, thick}
            ]

            % =========================
            % Step labels
            % =========================
            \node[anchor=east] at (-1.0,2.6) {\textbf{K-step}};
            \node[anchor=east] at (-1.0,1.25) {\textbf{F-step}};
            \node[anchor=east] at (-1.0,-0.25) {\textbf{C-step}};

            % =========================
            % K-step
            % =========================
            \node[box] (b1) at (0.8,2.6) {$\mathcal{B}_1$};
            \node[box] (b2) at (3.6,2.6) {$\mathcal{B}_2$};
            \node[box] (bm) at (6.6,2.6) {$\mathcal{B}_M$};
            \draw[dashline] (b2.east) -- (bm.west);

            \node[note] at (7.85,2.6) {
                Identifies hidden local structures in heterogeneous data.
            };

            % =========================
            % F-step
            % =========================
            \node[widebox] (m1) at (0.8,1.25) {$\mathcal{M}_{1,1}, \ldots, \mathcal{M}_{1,K}$};
            \node[widebox] (m2) at (3.6,1.25) {$\mathcal{M}_{2,1}, \ldots, \mathcal{M}_{2,K}$};
            \node[widebox] (mm) at (6.6,1.25) {$\mathcal{M}_{M,1}, \ldots, \mathcal{M}_{M,K}$};
            \draw[dashline] (m2.east) -- (mm.west);

            \node[note] at (7.85,1.25) {
                Trains supervised models for each discovered local group.
            };

            % =========================
            % C-step
            % =========================
            \node[aggbox] (agg) at (3.6,-0.25) {COBRA-Based\\Aggregation};

            \node[note] at (7.85,-0.25) {
                Combines local model predictions into the final output.
            };

            % =========================
            % Arrows
            % =========================
            \draw[arrow] (b1.south) -- (m1.north);
            \draw[arrow] (b2.south) -- (m2.north);
            \draw[arrow] (bm.south) -- (mm.north);

            \draw[arrow] (m1.south east) -- (agg.north);
            \draw[arrow] (m2.south) -- (agg.north);
            \draw[arrow] (mm.south west) -- (agg.north);

        \end{tikzpicture}
    }

    \vspace{0.45cm}

    \begin{block}{Key idea}
        KFCProcedure provides the clusterwise workflow, while COBRA provides the reusable combiner for the C-step.
    \end{block}

\end{frame}

\begin{frame}[t]{Background: GradientCOBRA Algorithm Workflow}
    \centering
    \scriptsize
    \vspace{-0.20cm}

    \resizebox{0.96\textwidth}{!}{%
        \begin{tikzpicture}[
                box/.style={
                        rectangle,
                        rounded corners=6pt,
                        thick,
                        align=center,
                        minimum width=2.45cm,
                        minimum height=0.78cm,
                        font=\scriptsize
                    },
                data/.style={
                        box,
                        draw=blue!70,
                        fill=blue!8
                    },
                train/.style={
                        box,
                        draw=orange!85!black,
                        fill=orange!10
                    },
                agg/.style={
                        box,
                        draw=green!55!black,
                        fill=green!10
                    },
                output/.style={
                        box,
                        draw=purple!70,
                        fill=purple!8
                    },
                arrow/.style={
                        -{Latex[length=1.8mm]},
                        thick
                    },
            ]

            % =========================
            % Input
            % =========================
            \node[data] (data) at (0,0) {
                Input\\Data
            };

            % =========================
            % Training row
            % =========================
            \node[train] (trainset) at (3.4,1.35) {
                Training\\Set
            };

            \node[train] (estimator) at (6.6,1.35) {
                Estimators\\
                $M_1,\ldots,M_K$
            };

            % =========================
            % Aggregation row
            % =========================
            \node[agg] (aggset) at (3.4,-1.35) {
                Aggregation\\Set
            };

            \node[agg] (predspace) at (6.6,-1.35) {
                Prediction\\Space
            };

            \node[agg] (similarity) at (9.8,-1.35) {
                Similarity\\Measure
            };

            \node[agg] (optimize) at (13.0,-1.35) {
                Optimize\\Weight
            };

            \node[agg] (aggregate) at (16.2,-1.35) {
                COBRA\\Aggregation
            };

            \node[output] (final) at (19.4,-1.35) {
                Final\\Prediction
            };

            % =========================
            % Split from input data
            % =========================
            \draw[arrow] (data) -- (trainset);
            \draw[arrow] (data) -- (aggset);

            % =========================
            % Training flow
            % =========================
            \draw[arrow] (trainset) -- (estimator);

            % =========================
            % Estimators feed prediction space
            % =========================
            \draw[arrow] (estimator) -- (predspace);

            % =========================
            % Aggregation flow
            % =========================
            \draw[arrow] (aggset) -- (predspace);
            \draw[arrow] (predspace) -- (similarity);
            \draw[arrow] (similarity) -- (optimize);
            \draw[arrow] (optimize) -- (aggregate);
            \draw[arrow] (aggregate) -- (final);

        \end{tikzpicture}
    }

    \vspace{0.30cm}

    \begin{block}{Key idea}
        GradientCOBRA uses trained estimators to create a prediction space, then measures similarity and optimizes aggregation to produce the final prediction.
    \end{block}

    \vspace{-0.05cm}
    {\tiny Source: Adapted from \cite{has2023gradientcobra}, \textit{Gradient COBRA: A kernel-based consensual aggregation for regression}.}

\end{frame}


\subsection{Problem Statement}
\begin{frame}[t]{Problem Statement}
    \small

    \vspace{0.2cm}
    \textbf{Main Problem : }

    \vspace{0.25cm}

    \begin{enumerate}

        \item How can a machine learning package support clusterwise learning?

        \item How can multiple aggregation strategies be unified in one framework?
    \end{enumerate}

    \vspace{0.35cm}
    \begin{block}{Project Focus}
        KFCProcedure provides a modular structure for local learning,
        prediction aggregation, and extensible model development.
    \end{block}
\end{frame}

\subsection{Objective}

\begin{frame}[t]{Objective}
    \footnotesize

    \begin{itemize}
        \setlength\itemsep{0.24cm}

        \item To study and implement the theoretical foundations of
              \textbf{KFCProcedure} and \textbf{COBRA-based aggregation methods}.

        \item To develop \textbf{KFCProcedure} as a clusterwise learning framework
              based on the K-step, F-step, and C-step workflow.

        \item To redesign \textbf{GradientCOBRA} as a reusable combiner architecture
              that supports flexible aggregation strategies for the C-step.

        \item To build a modular Python package with clear documentation, practical examples,
              and a consistent API for researchers and practitioners.
    \end{itemize}

\end{frame}


